\Chapter{47}

\Verse{1}
{
    \zA{话说王夫人听见邢夫人来了，连忙迎着出去。 邢夫人犹不知贾母已知鸳鸯之 事，正还又来打听信息，进了院门，早有几个婆子悄悄的回了他，他才知道。 待要
        回去，里面已知； 又见王夫人接出来了，少不得进来。 先与贾母请安，贾母一声儿 不言语，自己也觉得愧悔。 凤姐儿早指一事回避了。 鸳鸯也自回房去生气。
        薛姨妈 王夫人等恐碍着邢夫人的脸面，也都渐渐退了。 邢夫人且不敢出去。 贾母见无人， 方说道：“我听见你替你老爷说媒来了。}
}
{
    \eA{As soon as Madame Wang, so runs our narrative, heard of Madame Hsing's
        arrival, she quickly went out to welcome her. Madame Hsing was not yet aware
        that dowager lady Chia had learnt everything connected with Yüan Yang's
        affair, and she was coming again to see which way the wind blew.}
}
{
}

\Verse{2}
{
    \zA{你倒也‘三从四德’的，只是这贤惠也太 过了!你们如今也是孙子儿子满眼了，你还怕他使性子。 我听见你还由着你老爷的 那性子闹。”
        邢夫人满面通红，回道：“我劝过几次不依。 老太太还有什么不知道 的呢？ 我也是不得已儿。” 贾母道：“他逼着你杀人，你也杀去？ 如今你也想想：你
        兄弟媳妇，本来老实，又生的多病多痛，上上下下，那不是他操心？ 你一个媳妇， 虽然帮着，也是天天‘丢下耙儿弄扫帚’。 凡百事情，我如今自己减了。}
}
{
    \eA{The moment, however, she stepped inside the courtyard-entrance, several
        matrons promptly explained to her, quite confidentially, that their old
        mistress had been told all only a few minutes back, and she meant to retrace
        her steps, (but she saw that) every inmate in the suite of rooms was already
        conscious of her presence.}
}
{
}

\Verse{3}
{
    \zA{他们两个 就有些不到的去处，有鸳鸯那孩子还心细些，我的事情，他还想着一点子：该要的， 他就要了来； 该添什么，他就趁空儿告诉他们添了。
        鸳鸯再不这么着，娘儿两个， 里头外头大的小的，那里不忽略一件半件？ 我如今反倒自己操心去不成？ 还是天天盘 算和他们要东要西去？
        我这屋里有的没有的剩了他一个，年纪也大些，我凡做事的 脾气性格儿，他还知道些。 他二则也还投主子的缘法，他也并不指着我和那位太太
        要衣裳去，又和那位奶奶要银子去。}
}
{
    \eA{When she caught sight, besides, of Madame Wang walking out to meet her she
        had no option but to enter. First and foremost, she paid her respects to
        dowager lady Chia, but old lady Chia did not address her a single remark, so
        she felt within herself smitten with shame and remorse. Lady Feng soon gave
        something or other as an excuse and withdrew.}
}
{
}

\Verse{4}
{
    \zA{所以这几年，一应事情，他说什么，从你小婶 和你媳妇起，至家下大大小小，没有不信的。 所以不单我得靠，连你小婶、媳妇也 都省心。
        我有了这么个人，就是媳妇、孙子媳妇想不到的，我也不得缺了，也没气 可生了。 这会子他去了，你们又弄什么人来我使？ 你们就弄他那么个真珠儿似的人
        来，不会说话也无用。}
}
{
    \eA{Yüan Yang then returned also quite alone to her chamber to give vent to her
        resentment; and Mrs. Hsüeh, Madame Wang and the other inmates, one by one,
        retired in like manner, for fear of putting Madame Hsing out of countenance.
        Madame Hsing, however, could not muster courage to beat a retreat. Dowager
        lady Chia noticed that there was no one but themselves in her apartments.}
}
{
}

\Verse{5}
{
    \zA{我正要打发人和你老爷说去，他要什么人，我这里有钱，叫 他只管一万八千的买去就是，要这个丫头，不能!留下他伏侍我几年，就和他日夜 伏侍我尽了孝的一样。
        你来的也巧，就去说，更妥当了。” 说毕，命人来：“请了姨太太你姑娘们来。 才高兴说个话儿，怎么又都散了！” 丫头忙答应找去了。 众人赶忙的又来。
        只有薛姨妈向那丫鬟道：“我才来了，又做 什么去？ 你就说我睡了。” 那丫头道：“好亲亲的姨太太，姨祖宗!我们老太太生气 呢。}
}
{
    \eA{"I hear," she remarked, "that you had come to play the part of a go-between
        for your lord and master! You can very well observe the three obediences and
        four virtues, but this softness of yours is a work of supererogation! You
        people have also got now a whole lot of grandchildren and sons. Do you still
        live in fear and trembling lest he should put his monkey up?}
}
{
}

\Verse{6}
{
    \zA{你老人家不去，没个开交了。 只当疼我们罢!你老人家怕走，我背了你老人家 去。” 薛姨妈笑道：“小鬼头儿!你怕什么？ 不过骂几句就完了。” 说着，只得和这
        小丫头子走来。 贾母忙让坐，又笑道：“咱们斗牌罢？ 姨太太的牌也生了，咱们一 处坐着，别叫凤丫头混了我们去。”
        薛姨妈笑道：“正是呢，老太太替我看着些儿。 就是咱们娘儿四个斗呢，还是添一两个人呢？” 王夫人笑道：“可不只四个人？”
        凤姐儿道：“再添一个人，热闹些。”}
}
{
    \eA{Rumour has it that you yet let that disposition of your husband's run riot!"
        Madame Hsing's whole face got suffused with blushes. "I advised him time and
        again," she explained, "but he wouldn't listen to me. How is it, venerable
        senior, that you don't yet know that he turns a deaf ear to me? That's why I
        had no choice in the matter!"}
}
{
}

\Verse{7}
{
    \zA{贾母道：“叫鸳鸯来，叫他在这下手里坐着。 姨太太的眼花了，咱们两个的牌，都叫他看着些儿。” 凤姐笑了一声，向探春道： “你们知书识字的，倒不学算命？”
        探春道：“这又奇了，这会子你不打点精神赢 老太太几个钱，又想算命？” 凤姐儿道：“我正要算算今儿该输多少。 我还想赢呢？
        你瞧瞧，场儿没上，左右都埋伏下了。” 说的贾母薛姨妈都笑起来。 一时鸳鸯来了，便坐在贾母下首。 鸳鸯之下，便是凤姐儿。 铺下红毡，洗牌告
        么，五人起牌，斗了一回。}
}
{
    \eA{"Would you go and kill any one," dowager lady Chia asked, "that he might
        instigate you to? But consider now. Your brother's wife is naturally a quiet
        sort of person, and is born with many ailments; but is there anything,
        whether large or small, that she doesn't go to the trouble of looking after?}
}
{
}

\Verse{8}
{
    \zA{鸳鸯见贾母的牌已十成，只等一张二饼，便递了暗号儿 与凤姐儿。 凤姐儿正该发牌，便故意踌躇了半晌，笑道：“我这一张牌定在姨妈手
        里扣着呢，我若不发这一张牌，再顶不下来的。” 薛姨妈道：“我手里并没有你的 牌。” 凤姐儿道：“我回来是要查的。” 薛姨妈道：“你只管查。
        你且发下来，我 瞧瞧是张什么。” 凤姐儿便送在薛姨妈跟前，薛姨妈一看，是个二饼，便笑道：“我 倒不稀罕他，只怕老太太满了。”
        凤姐听了，忙笑道：“我发错了！”}
}
{
    \eA{And notwithstanding that that daughter-in-law of yours lends her a helping
        hand, she is daily so busy that she 'no sooner puts down the pick than she
        has to take up the broom.' So busy, that I have myself now curtailed a
        hundred and one things. But whenever there's anything those two can't
        manage, there's Yüan Yang to come to their assistance. She is, it's true, a
        mere child, but nevertheless very careful;}
}
{
}

\Verse{9}
{
    \zA{贾母笑的已 掷下牌来，说：“你敢拿回去!谁叫你错的不成？” 凤姐儿道：“可是我要算一算 命呢。 这是自己发的，也怨不得人了。”
        贾母笑道：“可是你自己打着你那嘴，问 着你自己才是。” 又向薛姨妈笑道：“我不是小气爱赢钱，原是个彩头儿。” 薛姨 妈笑道：“我们可不是这样想？
        那里有那样糊涂人，说老太太爱钱呢？” 凤姐儿正 数着钱，听了这话，忙又把钱穿上了，向众人笑道：“够了我的了!竟不为赢钱， 单为赢彩头儿。}
}
{
    \eA{and knows how to concern herself about my affairs a bit; indenting for
        anything that need be indented, and availing herself of an opportunity to
        tell them to supply every requisite. Were Yüan Yang not the kind of girl she
        is, how could those two ladies not neglect a whole or part of those matters,
        both important as well as unimportant, connected with the inner and outer
        quarters?}
}
{
}

\Verse{10}
{
    \zA{我到底小气，输了就数钱，快收起来罢。” 贾母规矩是鸳鸯代洗牌 的，便和薛姨妈说笑。 不见鸳鸯动手。 贾母道：“你怎么恼了，连牌也不替我洗？”
        鸳鸯拿起牌来笑道：“奶奶不给钱么！” 贾母道：“他不给钱，那是他交运了！” 便命小丫头子：“把他那一吊钱都拿过来！” 小丫头子真就拿了，搁在贾母傍边。
        凤姐儿笑道：“赏我罢，照数儿给就是了。” 薛姨妈笑道：“果然凤姐儿小气，不 过玩儿罢了。”}
}
{
    \eA{Would I not at present have to worry my own mind, instead of leaving things
        to others? Why, I'd daily have to rack my brain and go and ask them to give
        me whatever I might need! Of those girls, who've come to my quarters and
        those who've gone, there only remains this single one. She's, besides other
        respects, somewhat older in years, and has as well a slight conception of my
        ways of doing things, and of my tastes.}
}
{
}

\Verse{11}
{
    \zA{凤姐儿听说便站起来拉住薛姨妈，回头指着贾母素日放钱的一个木 箱子笑道：“姑妈瞧瞧，那个里头不知玩了我多少去了。 这一吊钱玩不了半个时辰，
        那里头的钱就招手儿叫他了。 只等把这一吊也叫进去了，牌也不用斗了，老祖宗气 也平了，又有正经事差我办去了。” 话未说完，引的贾母众人笑个不住。 正说着，
        偏平儿怕钱不够，又送了一吊来。 凤姐儿道：“不用放在我跟前，也放在老太太的 那一处罢。 一齐叫进去倒省事，不用做两次，叫箱子里的钱费事。”}
}
{
    \eA{In the second place, she has managed to win her mistresses' hearts, for she
        never tries to extort aught from me, or to dun this lady for clothes or that
        one for money. Hence it is that beginning from your sister-in-law and
        daughter-in-law down to the servants in the house, irrespective of old or
        young, there isn't a soul, who doesn't readily believe every single word she
        says in anything, no matter what it is!}
}
{
}

\Verse{12}
{
    \zA{贾母笑的手里 的牌撒了一桌子，推着鸳鸯，叫：“快撕他的嘴！” 平儿依言放下钱，也笑了一回，方回来。 至院门前，遇见贾琏，问他：“太太 在那里呢？
        老爷叫我请过去呢。” 平儿忙笑道：“在老太太跟前站了这半日，还没 动呢。 趁早儿丢开手罢。 老太太生了半日气，这会子亏二奶奶凑了半日的趣儿，才
        略好了些。” 贾琏道：“我过去，只说讨老太太示下，十四往赖大家去不去，好预 备轿子。 又请了太太，又凑了趣儿，岂不好呢。”}
}
{
    \eA{Not only do I thus have some one upon whom I can rely, but your young
        sister-in-law and your daughter-in-law are both as well spared much trouble.
        With a person such as this by me, should even my daughter-in-law and
        granddaughter-in-law not have the time to think of anything, I am not left
        without it; nor am I given occasion to get my temper ruffled.}
}
{
}

\Verse{13}
{
    \zA{平儿笑道：“依我说，你竟别过 去罢。 合家子连太太宝玉都有了不是，这会子你又填限去了。” 贾琏道：“已经完 了，难道还找补不成？ 况且与我又无干。
        二则老爷亲自吩咐我请太太去，这会子我 打发了人去，倘或知道了，正没好气呢，指着这个拿我出气罢。” 说着就走。 平儿 见他说的有理，也就跟了贾琏过来。
        到了堂屋里，便把脚步放轻了，往里间探头， 只见邢夫人站在那里。 凤姐儿眼尖，先瞧见了，便使眼色儿，不命他进来，又使眼 色与邢夫人。}
}
{
    \eA{But were she now to go, what kind of creature would they hunt up again to
        press into my service? Were you even to bring me a person made of real
        pearls, she'd be of no use; if she doesn't know how to speak!}
}
{
}

\Verse{14}
{
    \zA{邢夫人不便就走，只得倒了一碗茶来，放在贾母跟前。 贾母一回身， 贾琏不防，便没躲过。 贾母便问：“外头是谁？ 倒像个小子一伸头的似的。” 凤姐
        儿忙起身说：“我也恍惚看见有一个人影儿。” 一面说，一面起身出来。 贾琏忙进 去，陪笑道：“打听老太太十四可出门？ 好预备轿子。”
        贾母道：“既这么样，怎 么不进来，又做神做鬼的？” 贾琏陪笑道：“见老太太玩牌，不敢惊动，不过叫媳 妇出来问问。”}
}
{
    \eA{I was just about to send some one to go and explain to your husband that
        'I've got money in here enough to buy any girl he fancies,' and to tell him
        that 'he's at liberty to give for her purchase from eight to ten thousand
        taels; that, if he has set his heart upon this girl, he can't however have
        her;}
}
{
}

\Verse{15}
{
    \zA{贾母道：“就忙到这一时!等他家去，你问他多少问不得？ 那一遭儿 你这么小心来？ 这又不知是来做耳报神的，也不知是来做探子的，鬼鬼祟祟，倒吓 我一跳。
        什么好下流种子!你媳妇和我玩牌呢，还有半日的空儿，你家去再和那赵 二家的商量治你媳妇去罢！” 说着众人都笑了。 鸳鸯笑道：“鲍二家的，老祖宗又
        拉上赵二家的去。” 贾母也笑道：“可不？ 我那里记得什么‘抱’着‘背’着的。 提起这些事来，不由我不生气。}
}
{
    \eA{and that by leaving her behind to attend to me, during the few years to
        come, it will be just the same as if he tried to acquit himself of his
        filial duties by waiting upon me day and night,' so you come at a very
        opportune moment. Were you therefore to go yourself at once and deliver him
        my message, it will answer the purpose far better!" These words over, she
        called the servants.}
}
{
}

\Verse{16}
{
    \zA{我进了这门子做重孙媳妇起，到如今我也有个重孙 子媳妇了，连头带尾五十四年，凭着大惊大险、千奇百怪的事也经了些，从没经过 这些事。 还不离了我这里呢！”
        贾琏一声儿不敢说，忙退出来。 平儿在窗外站着，悄悄的笑道：“我说你不听， 到底碰在网里了。” 正说着，只见邢夫人也出来。 贾琏道：“都是老爷闹的，如今
        都搁在我和太太身上。” 邢夫人道：“我把你这没孝心的种子!人家还替老子死呢。 白说了几句，你就抱怨天、抱怨地了。}
}
{
    \eA{"Go," she said, "and ask Mrs. Hsüeh, and your young mistresses to come! We
        were in the middle of a chat full of zest, and how is it they've all
        dispersed?" The waiting-maids immediately assented and left to go in search
        of their mistresses, one and all of whom promptly re-entered her apartments,
        with the sole exception of Mrs. Hsüeh.}
}
{
}

\Verse{17}
{
    \zA{你还不好好的呢!这几日生气，仔细他捶你。” 贾琏道：“太太快过去罢，叫我来请了好半日了。” 说着，送他母亲出来过那边去。
        邢夫人将方才的话只略说了几句，贾赦无法，又且含愧，自此便告了病，且不 敢见贾母，只打发邢夫人及贾琏每日过去请安。 只得又各处遣人购求寻觅，终久费
        了五百两银子买了一个十七岁女孩子来，名唤嫣红，收在屋里，不在话下。 这里斗 了半日牌，吃晚饭才罢。 此一二日间无话。
        转眼到了十四，黑早，赖大的媳妇又进来请。}
}
{
    \eA{"I've only now returned," she observed to the waiting-maid, "and what shall
        I go again for? Just tell her that I'm fast asleep!" "Dearest Mrs. Hsüeh!"
        the waiting-maid pleaded, "my worthy senior! our old mistress will get
        angry. If you, venerable lady, don't appear nothing will appease her; so do
        it for the love of us! Should you object to walking, why I'm quite ready to
        carry you on my back." "You little imp!"}
}
{
}

\Verse{18}
{
    \zA{贾母高兴，便带了王夫人薛姨妈 及宝玉姐妹等至赖大花园中，坐了半日。 那花园虽不及大观园，却也十分齐整宽阔， 泉石林木，楼台亭轩，也有好几处动人的。
        外面大厅上，薛蟠、贾珍、贾琏、贾蓉 并几个近族的都来了。 那赖大家内，也请了几个现任的官长并几个大家子弟作陪。
        因其中有个柳湘莲，薛蟠自上次会过一次，已念念不忘。}
}
{
    \eA{Mrs. Hsüeh laughed. "What are you afraid of? All she'll do will be to scold
        you a little; and it will all be over soon!" While replying, she felt that
        she had no course but to retrace her footsteps, in company with the
        waiting-maid. Dowager lady Chia at once motioned her into a seat. "Let's
        have a game of cards!" she then smilingly proposed. "You, Mrs. Hsüeh, are
        not a good hand at them;}
}
{
}

\Verse{19}
{
    \zA{又打听他最喜串戏，且都 串的是生旦风月戏文，不免错会了意，误认他做了风月子弟，正要与他相交，恨没 有个引进，这一天可巧遇见，乐得无可不可。
        且贾珍等也慕他的名，酒盖住了脸， 就求他串了两出戏。 下来，移席和他一处坐着，问长问短，说东说西。 那柳湘莲原
        系世家子弟，读书不成，父母早丧，素性爽侠，不拘细事，酷好耍枪舞剑，赌博吃 酒，以至眠花卧柳，吹笛弹筝，无所不为。 因他年纪又轻，生得又美，不知他身分
        的人，都误认作优伶一类。}
}
{
    \eA{so let's sit together, and see that lady Feng doesn't cheat us!" "Quite so,"
        laughed Mrs. Hsüeh. "But it will be well if your venerable ladyship would
        look over my hand a bit! Are we four ladies to play, or are we to add one or
        two more persons to our number?" "Naturally only four!" Madame Wang smiled.
        "Were one more player let in," lady Feng interposed, "it would be merrier!"}
}
{
}

\Verse{20}
{
    \zA{那赖大之子赖尚荣与他素昔交好，故今儿请来做陪。 不 想酒后别人犹可，独薛蟠又犯了旧病。 心中早已不快，得便意欲走开完事。 无奈赖
        尚荣又说：“方才宝二爷又嘱咐我：才一进门，虽见了，只是人多不好说话，叫我 嘱咐你散的时候别走，他还有话说呢。 你既一定要去，等我叫出他来，你两个见了
        再走，与我无干。” 说着，便命小厮们：“到里头，找一个老婆子，悄悄告诉，请 出宝二爷来。” 那小厮去了。 没一杯茶时候，果见宝玉出来了。}
}
{
    \eA{"Call Yüan Yang here," old lady Chia suggested, "and make her take this
        lower seat; for as Mrs. Hsüeh's eyesight is rather dim, we'll charge her to
        look over our two hands a bit." "You girls know how to read and write," lady
        Feng remarked with a smile, addressing herself to T'an Ch'un, "and why don't
        you learn fortune-telling?" "This is again strange!" T'an Ch'un exclaimed.}
}
{
}

\Verse{21}
{
    \zA{赖尚荣向宝玉笑道：“好叔叔，把他交给你， 我张罗人去了。” 说着已经去了。 宝玉便拉了柳湘莲到厅侧书房坐下，问他：“这 几日可到秦钟的坟上去了？”
        湘莲道：“怎么不去？ 前儿我们几个放鹰去，离他坟 上还有二里，我想今年夏天雨水勤，恐怕他坟上站不住。 我背着众人走到那里去瞧
        了一瞧，略又动了一点子，回家来就便弄了几百钱，第三日一早出去雇了两个人收 拾好了。” 宝玉说：“怪道呢。}
}
{
    \eA{"Instead of bracing up your energies now to rook some money out of our
        venerable senior, you turn your thoughts to fortune-telling!" "I was just
        wishing to consult the fates," lady Feng proceeded, "as to how much I shall
        lose to-day. Can I ever dream of winning? Why, look here. We haven't
        commenced playing, and they have placed themselves in ambush on the left and
        right."}
}
{
}

\Verse{22}
{
    \zA{上月我们大观园的池子里头结了莲蓬，我摘了十个， 叫焙茗出去到坟上供他去。 回来我也问他：‘可被雨冲坏了没有？’ 他说：‘不但 没冲，更比上回新了些。’
        我想着必是这几个朋友新收拾了。 我只恨我天天圈在家 里，一点儿做不得主，行动就有人知道，不是这个拦就是那个劝的，能说不能行。 虽然有钱，又不由我使。”
        柳湘莲道：“这个事也用不着你操心，外头有我，你只 心里有了就是了。 眼前十月初一日，我已经打点下上坟的花销。}
}
{
    \eA{This remark amused dowager lady Chia and Mrs. Hsüeh. But presently Yüan Yang
        arrived, and seated herself below her old mistress. After Yüan Yang sat lady
        Feng. The red cloth was then spread; the cards were shuffled; the dealer was
        decided upon and the quintet began to play.}
}
{
}

\Verse{23}
{
    \zA{你知道，我一贫如 洗，家里是没的积聚的； 纵有几个钱来，随手就光的。 不如趁空儿留下这一分，省 的到了跟前扎煞手。”
        宝玉道：“我也正为这个，要打发焙茗找你。 你又不大在家， 知道你天天萍踪浪迹，没个一定的去处。” 柳湘莲道：“你也不用找我，这个事也 不过各尽其道。
        眼前我还要出门去走走，外头游逛三年五载再回来。” 宝玉听了， 忙问：“这是为何？” 柳湘莲冷笑道：“我的心事，等到跟前，你自然知道。 我如
        今要别过了。”}
}
{
    \eA{After the game had gone on for a time, Yüan Yang noticed that dowager lady
        Chia had a full hand and was only waiting for one two-spotted card, and she
        made a secret sign to lady Feng. Lady Feng was about to lead, but purposely
        lingered for a few moments. "This card will, for a certainty, be snatched by
        Mrs. Hsüeh," she smiled, "yet if I don't play this one, I won't be able
        later to come out with what I want."}
}
{
}

\Verse{24}
{
    \zA{宝玉道：“好容易会着，晚上同散，岂不好？” 湘莲道：“你那令 姨表兄还是那样，再坐着未免有事，不如我回避了倒好。” 宝玉想一想，说道：“既
        是这么样，倒是回避他为是。 只是你要果真远行，必须先告诉我一声，千万别悄悄 的去了。” 说着，便滴下泪来。 柳湘莲说道：“自然要辞你去，你只别和别人说就
        是了。” 说着就站起来要走； 又道：“你就进去罢，不必送我。” 一面说，一面出了书房。 刚至大门前，早遇见薛蟠在那里乱叫：“谁放了小柳 儿走了？”}
}
{
    \eA{"I haven't got any cards you want in my hand," Mrs. Hsüeh remarked. "I mean
        to see by and bye," lady Feng resumed. "You're at liberty to see," Mrs.
        Hsüeh said. "But go on, play now! Let me look what card it is." Lady Feng
        threw the card in front of Mrs. Hsüeh. At a glance, Mrs. Hsüeh perceived
        that it was the two spot. "I don't fancy this card," she smiled. "What I
        fear is that our dear senior will get a full hand."}
}
{
}

\Verse{25}
{
    \zA{柳湘莲听了，火星乱迸，恨不得一拳打死； 复思酒后挥拳，又碍着赖尚 荣的脸面，只得忍了又忍。 薛蟠忽见他走出来，如得了珍宝，忙趔趄着，走上去一
        把拉住，笑道：“我的兄弟!你往那里去了？” 湘莲道：“走走就来。” 薛蟠笑道：
        “你一去都没了兴头了，好歹坐一坐，就算疼我了!凭你什么要紧的事，交给哥哥， 只别忙。 你有这个哥哥，你要做官发财都容易。”}
}
{
    \eA{"I've played wrong!" lady Feng laughingly exclaimed at these words. Dowager
        lady Chia laughed, and throwing down her cards, "If you dare," she shouted,
        "take it back! Who told you to play the wrong card?" "Didn't I want to have
        my fortune told?" lady Feng observed. "I played this card of my own accord,
        so there's no one with whom I can find fault." "You should then beat your
        own lips and punish your own self;}
}
{
}

\Verse{26}
{
    \zA{湘莲见他如此不堪，心中又恨又 恼，早生一计，拉他到僻净处，笑道：“你真心和我好，还是假心和我好呢？” 薛
        蟠听见这话，喜得心痒难挠，乜斜着眼，笑道：“好兄弟!你怎么问起我这样话来？ 我要是假心，立刻死在眼前。” 湘莲道：“既如此，这里不便。
        等坐一坐，我先走， 你随后出来，跟到我下处，咱们索性喝一夜酒。 我那里还有两个绝好的孩子，从没 出门的。
        你可连一个跟的人也不用带，到了那里，伏侍人都是现成的。”}
}
{
    \eA{it's only fair;" old lady Chia remarked. Then facing Mrs. Hsüeh, "I'm not a
        niggard, fond of winning money," she went on to say, "but it was my good
        luck!" "Don't we too think as much?" Mrs. Hsüeh smiled. "Who's there stupid
        enough to say that your venerable ladyship's heart is set upon money?" Lady
        Feng was busy counting the cash, but catching what was said, she restrung
        them without delay.}
}
{
}

\Verse{27}
{
    \zA{薛蟠听如 此说，喜的酒醒了一半，说：“果然如此？” 湘莲笑道：“如何!人拿真心待你， 你倒不信了。” 薛蟠忙笑道：“我又不是呆子，怎么有个不信的呢？
        既如此，我又 不认得，你先去了，我在那里找你？” 湘莲道：“我这下处在北门外头，你可舍得 家，城外住一夜去？” 薛蟠道：“有了你，我还要家做什么！”
        湘莲道：“既如此， 我在北门外头桥上等你。 咱们席上且吃酒去。 你看我走了之后你再走，他们就不留 神了。” 薛蟠听了，连忙答应道是。}
}
{
    \eA{"I've got my share," she said, laughingly to the company. "It isn't at all
        that you wish to win. It's your good luck that made you come out a winner!
        But as for me, I am really a mean creature; and, as I managed to lose, I
        count the money and put it away at once."}
}
{
}

\Verse{28}
{
    \zA{二人复又入席，饮了一回。 那薛蟠难熬，只拿眼看湘莲，心内越想越乐，左一 壶，右一壶，并不用人让，自己就吃了又吃，不觉酒有八九分了。 湘莲就起身出来，
        瞅人不防出至门外，命小厮杏奴：“先家去罢，我到城外就来。” 说毕，已跨马直 出北门，桥上等候薛蟠。 一顿饭的工夫，只见薛蟠骑着一匹马，远远的赶了来，张
        着嘴，瞪着眼，头似拨浪鼓一般，不住左右乱瞧。 及至从湘莲马前过去，只顾往远 处瞧，不曾留心近处。 湘莲又笑又恨，他便也撒马随后跟来。}
}
{
    \eA{Dowager lady Chia usually made Yüan Yang shuffle the cards for her, but
        being engaged in chatting and joking with Mrs. Hsüeh, she did not notice
        Yüan Yang take them in hand. "Why is it you're so huffed," old lady Chia
        asked, "that you don't even shuffle for me?" "Lady Feng won't let me have
        the money!" Yüan Yang replied, picking up the cards.}
}
{
}

\Verse{29}
{
    \zA{薛蟠往前看时，渐渐 人烟稀少，便又圈马回来，再不想一回头见了湘莲，如获奇珍，忙笑道：“我说你 是个再不失信的。”
        湘莲笑道：“快往前走，仔细人看见跟了来，就不好了。” 说 着，先就撒马前去。 薛蟠也就紧紧跟来。
        湘莲见前面人烟已稀，且有一带苇塘，便下马，将马拴在树上，向薛蟠笑道： “你下来，咱们先设个誓。 日后要变了心，告诉别人的，就应誓。” 薛蟠笑道：“这
        话有理。”}
}
{
    \eA{"If she doesn't give the money," dowager lady Chia observed, "it will be a
        turning-point in her luck. Take that string of a thousand cash of hers," she
        accordingly directed a servant, "and bring it bodily over here!" A young
        waiting-maid actually fetched the string of cash and deposited it by the
        side of her old mistress.}
}
{
}

\Verse{30}
{
    \zA{连忙下马，也拴在树上，便跪下说道：“我要日久变心，告诉人去的， 天诛地灭。” 一言未了，只听“镗”的一声，背后好似铁锤砸下来，只觉得一阵黑，
        满眼金星乱迸，身不由己，就倒在地下了。 湘莲走上来瞧瞧，知道他是个不惯捱打 的，只使了三分气力，向他脸上拍了几下，登时便开了果子铺。 薛蟠先还要扎挣起
        身，又被湘莲用脚尖点了一点，仍旧跌倒。 口内说道：“原来是两家情愿，你不依， 只管好说，为什么哄出我来打我？” 一面说，一面乱骂。}
}
{
    \eA{"Let me have them," lady Feng eagerly cried smiling, "and I'll square all
        that's due, and finish." "In very truth, lady Feng, you're a miserly
        creature!" Mrs. Hsüeh laughed. "It's simply for mere fun, nothing more!"
        Lady Feng, at this insinuation, speedily stood up, and, laying her hand on
        Mrs. Hsüeh, she turned her head round, and pointed at a large wooden box, in
        which old lady Chia usually deposited her money.}
}
{
}

\Verse{31}
{
    \zA{湘莲道：“我把你这瞎了 眼的，你认认柳大爷是谁!你不说哀求，你还伤我!我打死你也无益，只给你个利害 罢。”
        说着，便取了马鞭过来，从背后至胫，打了三四十下。 薛蟠的酒早已醒了大 半，不觉得疼痛难禁，由不的“嗳哟”一声。 湘莲冷笑道：“也只如此，我只当你
        是不怕打的。” 一面说，一面又把薛蟠的左腿拉起来，向苇中泞泥处拉了几步，滚 的满身泥水，又问道：“你可认得我了？” 薛蟠不应，只伏着哼哼。}
}
{
    \eA{"Aunt," she said, a smile curling her lips, "look here! I couldn't tell you
        how much there is in that box that was won from me! This tiao will be
        wheedled by the cash in it, before we've played for half an hour! All we've
        got to do is to give them sufficient time to lure this string in as well; we
        needn't trouble to touch the cards. Your temper, worthy ancestor, will thus
        calm down.}
}
{
}

\Verse{32}
{
    \zA{湘莲又掷下鞭 子，用拳头向他身上擂了几下，薛蟠便乱滚乱叫，说：“肋条折了!我知道你是正 经人，因为我错听了旁人的话了！”
        湘莲道：“不用拉旁人，你只说现在的。” 薛 蟠道：“现在也没什么说的，不过你是个正经人，我错了！” 湘莲道：“还要说软 些，才饶你。”
        薛蟠哼哼的道：“好兄弟――”湘莲便又一拳。 薛蟠“嗳”了一声 道：“好哥哥――”湘莲又连两拳。 薛蟠忙嗳哟叫道：“好老爷!饶了我这没眼睛
        的瞎子罢!从今以后，我敬你怕你了！”}
}
{
    \eA{If you've also got any legitimate thing for me to do, you might bid me go
        and attend to it!" This joke had scarcely been concluded than it evoked
        incessant laughter from dowager lady Chia and every one else. But while she
        was bandying words, P'ing Erh happened to bring her another string of cash
        prompted by the apprehension that her capital might not suffice to meet her
        wants.}
}
{
}

\Verse{33}
{
    \zA{湘莲道：“你把那水喝两口。” 薛蟠一面 听了，一面皱眉道：“这水实在腌？ ，怎么喝的下去！” 湘莲举拳就打，薛蟠忙道： “我喝我喝！”
        说着，只得俯头向苇根下喝了一口，犹未咽下去，只听“哇”的一 声，把方才吃的东西都吐了出来。 湘莲道：“好腌？ 东西，你快吃完了，饶你。”
        薛蟠听了，叩头不迭，说：“好歹积阴功饶我罢!这至死不能吃的。” 湘莲道：“这 么气息，倒熏坏了我！” 说着，丢下了薛蟠，便牵马认镫去了。}
}
{
    \eA{"It's useless putting them in front of me!" lady Feng cried. "Place these
        too over there by our old lady and let them be wheedled in along with the
        others! It will thus save trouble, as there won't be any need to make two
        jobs of them, to the inconvenience of the cash already in the box." Dowager
        lady Chia had a hearty laugh, so much so, that the cards, she held in her
        hand, flew all over the table;}
}
{
}

\Verse{34}
{
    \zA{这里薛蟠见他已去， 方放下心来，后悔自己不该误认了人。 待要扎挣起来，无奈遍体疼痛难禁。 谁知贾珍等席上忽不见了他两个，各处寻找不见。
        有人说：“恍惚出北门去了。” 薛蟠的小厮素日是惧他的，他吩咐了不许跟去，谁敢找去。 后来还是贾珍不放心，
        命贾蓉带着小厮们寻踪问迹的，直找出北门，下桥二里多路，忽见苇坑旁边薛蟠的 马拴在那里。 众人都道：“好了，有马必有人。” 一齐来至马前，只听苇中有人呻
        吟。}
}
{
    \eA{but pushing Yüan Yang. "Be quick," she shouted, "and wrench that mouth of
        hers!" P'ing Erh placed the cash according to her mistress' directions. But
        after indulging too in laughter for a time, she retraced her footsteps. On
        reaching the entrance into the court, she met Chia Lien. "Where's your
        Madame Hsing?" he inquired. "Mr. Chia She told me to ask her to go round."}
}
{
}

\Verse{35}
{
    \zA{大家忙走来一看，只见薛蟠的衣衫零碎，面目肿破，没头没脸，遍身内外滚的 似个泥母猪一般。 贾蓉心内已猜着八九了，忙下马命人搀了起来，笑道：“薛大叔
        天天调情，今日调到苇子坑里。 必定是龙王爷也爱上你风流，要你招驸马去，你就 碰到龙犄角上了！” 薛蟠羞的没地缝儿钻进去，那里爬的上马去？ 贾蓉命人赶到关
        厢里雇了一乘小轿子，薛蟠坐了，一齐进城。 贾蓉还要抬往赖家去赴席，薛蟠百般 苦告，央及他不用告诉人，贾蓉方依允了，让他各自回家。}
}
{
    \eA{"She's been standing in there with our old mistress," P'ing Erh hastily
        laughed, "for ever so long, and yet she isn't inclined to budge! Seize the
        earliest opportunity you can get to wash your hands clean of this business!
        Our old lady has had a good long fit of fuming and raging. Luckily, our lady
        Secunda cracked an endless stock of jokes, so she, at length, got a bit
        calmer!" "I'll go over," Chia Lien said.}
}
{
}

\Verse{36}
{
    \zA{贾蓉仍往赖家回复贾珍 并方才的形景。 贾珍也知湘莲所打，也笑道：“他须得吃个亏才好。” 至晚散了， 便来问候。 薛蟠自在卧房将养，推病不见。
        贾母等回来各自归家时，薛姨妈与宝钗见香菱哭的眼睛肿了，问起原故，忙来 瞧薛蟠时，脸上身上虽见伤痕，并未伤筋动骨。 薛姨妈又是心疼，又是发恨，骂一
        回薛蟠，又骂一回湘莲，意欲告诉王夫人，遣人寻拿湘莲。 宝钗忙劝道：“这不是 什么大事，不过他们一处吃酒，酒后反脸常情。 谁醉了，多挨几下子打，也是有的。}
}
{
    \eA{"All I have to do is to find out our venerable senior's wishes, as to
        whether she means to go to Lai Ta's house on the fourteenth, so that I might
        have time to get the chairs ready. As I'll be able to tell Madame Hsing to
        return, and have a share of the fun, won't it be well for me to go?" "My
        idea is," P'ing Erh suggested laughingly, "that you shouldn't put your foot
        in there!}
}
{
}

\Verse{37}
{
    \zA{况且咱们家的无法无天的人，也是人所共知的。 妈妈不过是心疼的原故，要出气也 容易。 等三五天哥哥好了出得去的时候，那边珍大爷琏二爷这干人也未必白丢开
        手，自然备个东道，叫了那个人来，当着众人替哥哥赔不是认罪就是了。 如今妈妈 先当件大事告诉众人，倒显的妈妈偏心溺爱，纵容他生事招人，今儿偶然吃了一次
        亏，妈妈就这样兴师动众，倚着亲戚之势欺压常人。” 薛姨妈听了道：“我的儿! 到底是你想的到，我一时气糊涂了。” 宝钗笑道：“这才好呢。}
}
{
    \eA{Every one, even up to Madame Wang, and Pao-yü, have alike received a rap on
        the knuckles, and are you also going now to fill up the gap?" "Everything is
        over long ago," Chia Lien observed, "and can it be that she'll cap the whole
        thing by blowing me up too? What's more, it's no concern of mine.}
}
{
}

\Verse{38}
{
    \zA{他又不怕妈妈，又 不听人劝，一天纵似一天。 吃过两三个亏，他也罢了。” 薛蟠睡在炕上，痛骂湘莲，又命小厮：“去拆他的房子，打死他，和他打官司！”
        薛姨妈喝住小厮们，只说：“湘莲一时酒后放肆，如今酒醒，后悔不及，惧罪逃走 了。” 薛蟠听见如此说了，要知端底，且看下回分解。 (本章完)}
}
{
    \eA{In the next place, Mr. Chia She enjoined me that I was to go in person, and
        ask his wife round, so, if I at present depute some one else, and he comes
        to know about it, he really won't feel in a pleasant mood, and he'll take
        advantage of this pretext to give vent to his spite on me." These words
        over, he quickly marched off.}
}
{
}

\Verse{39}
{
    \zA{}
}
{
    \eA{And P'ing Erh was so impressed with the reasonableness of his arguments,
        that she followed in his track. As soon as Chia Lien reached the reception
        hall, he trod with a light step. Then peeping in he saw Madame Hsing
        standing inside. Lady Feng, with her eagle eye, was the first to espy him.
        But she winked at him and dissuaded him from coming in, and next gave a wink
        to Madame Hsing. Madame Hsing could not conveniently get away at once, and
        she had to pour a cup of tea, and place it in front of dowager lady Chia.
        But old lady Chia jerked suddenly round, and took Chia Lien at such a
        disadvantage that he found it difficult to beat a retreat. "Who is outside?"
        exclaimed old lady Chia. "It seemed to me as if some servant-boy had poked
        his head in." Lady Feng sprung to her feet without delay.}
}
{
}

\Verse{40}
{
    \zA{}
}
{
    \eA{"I also," she interposed, "indistinctly noticed the shadow of some one."
        Saying this, she walked away and quitted the room. Chia Lien entered with
        hasty step. Forcing a smile, "I wanted to ask," he remarked, "whether you,
        venerable senior, are going out on the fourteenth, so that the chairs may be
        got ready." "In that case," dowager lady Chia rejoined, "why didn't you come
        straight in; but behaved again in that mysterious way?" "I saw that you were
        playing at cards, dear ancestor," Chia Lien explained with a strained laugh,
        "and I didn't venture to come and disturb you. I therefore simply meant to
        call my wife out to find out from her." "Is it anything so very urgent that
        you had to say it this very moment?" old lady Chia continued.}
}
{
}

\Verse{41}
{
    \zA{}
}
{
    \eA{"Had you waited until she had gone home, couldn't you have asked her any
        amount of questions you may have liked? When have you been so full of zeal
        before? I'm puzzled to know whether it isn't as an eavesdropping spirit that
        you appear on the scene; nor can I say whether you don't come as a spy. But
        that impish way of yours gave me quite a start! What a low-bred fellow you
        are! Your wife will play at cards with me for a good long while more, so
        you'd better bundle yourself home, and conspire again with Chao Erh's wife
        how to do away with your better half." Her remarks evoked general merriment.
        "It's Pao Erh's wife," Yüan Yang put in laughingly, "and you, worthy senior,
        have dragged in again Chao Erh's wife." "Yes!" assented old lady Chia,
        likewise with a laugh.}
}
{
}

\Verse{42}
{
    \zA{}
}
{
    \eA{"How could I remember whether he wasn't (pao) embracing her, or (pei)
        carrying her on his back. The bare mention of these things makes me lose all
        self-control and provokes me to anger! Ever since I crossed these doors as a
        great grandson's wife, I have never, during the whole of these fifty-four
        years, seen anything like these affairs, albeit it has been my share to go
        through great frights, great dangers, thousands of strange things and
        hundred and one remarkable occurrences! Don't you yet pack yourself off from
        my presence?"}
}
{
}

\Verse{43}
{
    \zA{}
}
{
    \eA{Chia Lien could not muster courage to utter a single word to vindicate
        himself, but retired out of the room with all promptitude. P'ing Erh was
        standing outside the window. "I gave you due warning in a gentle tone, but
        you wouldn't hear; you've, after all, rushed into the very meshes of the
        net!" These reproaches were still being heaped on him when he caught sight
        of Madame Hsing, as she likewise made her appearance outside. "My father,"
        Chia Lien ventured, "is at the bottom of all this trouble; and the whole
        blame now is shoved upon your shoulders as well as mine, mother." "I'll take
        you, you unfilial thing and…" Madame Hsing shouted. "People lay down their
        lives for their fathers;}
}
{
}

\Verse{44}
{
    \zA{}
}
{
    \eA{and you are prompted by a few harmless remarks to murmur against heaven and
        grumble against earth! Won't you behave in a proper manner? He's in high
        dudgeon these last few days, so mind he doesn't give you a pounding!"
        "Mother, cross over at once," Chia Lien urged; "for he told me to come and
        ask you to go a long time ago." Pressing his mother, he escorted her outside
        as far as the other part of the mansion. Madame Hsing gave (her husband)
        nothing beyond a general outline of all that had been recently said; but
        Chia She found himself deprived of the means of furthering his ends. Indeed,
        so stricken was he with shame that from that date he pleaded illness.}
}
{
}

\Verse{45}
{
    \zA{}
}
{
    \eA{And so little able was he to rally sufficient pluck to face old lady Chia,
        that he merely commissioned Madame Hsing and Chia Lien to go daily and pay
        their respects to her on his behalf. He had no help too but to despatch
        servants all over the place to make every possible search and inquiry for a
        suitable concubine for him. After a long time they succeeded in purchasing,
        for the sum of eighty taels, a girl of seventeen years of age, Yen Hung by
        name, whom he introduced as secondary wife into his household. But enough of
        this subject. In the rooms on the near side, they protracted for a long time
        their noisy game of cards, and only broke up after they had something to
        eat.}
}
{
}

\Verse{46}
{
    \zA{}
}
{
    \eA{Nothing worthy of note, however, occurred during the course of the following
        day or two. In a twinkle, the fourteenth drew near. At an early hour before
        daybreak, Lai Ta's wife came again into the mansion to invite her guests.
        Dowager lady Chia was in buoyant spirits, so taking along Madame Wang, Mrs.
        Hsüeh, Pao-yü and the various young ladies, she betook herself into Lai Ta's
        garden, where she sat for a considerable time. This garden was not, it is
        true, to be compared with the garden of Broad Vista; but it also was most
        beautifully laid out, and consisted of spacious grounds.}
}
{
}

\Verse{47}
{
    \zA{}
}
{
    \eA{In the way of springs, rockeries, arbours and woods, towers and terraces,
        pavilions and halls, it likewise contained a good many sufficient to excite
        admiration. In the main hall outside, were assembled Hsüeh P'an, Chia Chen,
        Chia Lien, Chia Jung and several close relatives. But Lai Ta had invited as
        well a number of officials, still in active service, and numerous young men
        of wealthy families, to keep them company. Among that party figured one Liu
        Hsiang-lien, whom Hsüeh P'an had met on a previous occasion and kept ever
        since in constant remembrance.}
}
{
}

\Verse{48}
{
    \zA{}
}
{
    \eA{Having besides discovered that he had a passionate liking for theatricals,
        and that the parts he generally filled were those of a young man or lady, in
        fast plays, he had unavoidably misunderstood the object with which he
        indulged in these amusements, to such a degree as to misjudge him for a
        young rake. About this time, he had been entertaining a wish to cultivate
        intimate relations with him, but he had, much to his disgust, found no one
        to introduce him, so when he, by a strange coincidence, came to be thrown in
        his way, on the present occasion, he revelled in intense delight.}
}
{
}

\Verse{49}
{
    \zA{}
}
{
    \eA{But Chia Chen and the other guests had heard of his reputation, so as soon
        as wine had blinded their sense of shame, they entreated him to sing two
        short plays; and when subsequently they got up from the banquet, they
        ensconced themselves near him, and, pressing him with questions, they
        carried on a conversation on one thing and then another. This Liu
        Hsiang-lien was, in fact, a young man of an old family; but he had been
        unsuccessful in his studies, and had lost his father and mother. He was
        naturally light-hearted and magnanimous; not particular in minor matters;
        immoderately fond of spear-exercise and fencing, of gambling and boozing;
        even going to such excesses as spending his nights in houses of easy virtue;
        playing the fife, thrumming the harp, and going in for everything and
        anything.}
}
{
}

\Verse{50}
{
    \zA{}
}
{
    \eA{Being besides young in years, and of handsome appearance, those who did not
        know what his standing was, invariably mistook him for an actor. But Lai
        Ta's son had all along been on such friendly terms with him, that he
        consequently invited him for the nonce to help him do the honours. Of a
        sudden, while every one was, after the wines had gone round, still on his
        good behaviour, Hsüeh P'an alone got another fit of his old mania.}
}
{
}

\Verse{51}
{
    \zA{}
}
{
    \eA{From an early stage, his spirits sunk within him and he would fain have
        seized the first convenient moment to withdraw and consummate his designs
        but for Lai Shang-jung, who then said: "Our Mr. Pao-yü told me again just
        now that although he saw you, as he walked in, he couldn't speak to you with
        so many people present, so he bade me ask you not to go, when the party
        breaks up, as he has something more to tell you. But as you insist upon
        taking your leave, you'd better wait until I call him out, and when you've
        seen each other, you can get away; I'll have nothing to say then." While
        delivering the message, "Go inside," he directed the servant-boys, "and get
        hold of some old matron and tell her quietly to invite Mr. Pao-yü to come
        out."}
}
{
}

\Verse{52}
{
    \zA{}
}
{
    \eA{A servant-lad went on the errand, and scarcely had time enough elapsed to
        enable one to have a cup of tea in, than Pao-yü, actually, made his
        appearance outside. "My dear sir," Lai Shang-jung smilingly observed to
        Pao-yü, "I hand him over to you. I'm going to entertain the guests!" With
        these words, he was off. Pao-yü pulled Lia Hsiang-lien into a side study in
        the hall, where they sat down. "Have you been recently to Ch'in Ch'ung's
        grave?" he inquired of him. "How could I not go?" Hsiang-lien answered. "The
        other day a few of us went out to give our falcons a fly; and we were yet at
        a distance of two li from his tomb, when remembering the heavy rains, we've
        had this summer, I gave way to fears lest his grave may not have been proof
        against them;}
}
{
}

\Verse{53}
{
    \zA{}
}
{
    \eA{so evading the notice of the party I went over and had a look. I found it
        again slightly damaged; but when I got back home, I speedily raised a few
        hundreds of cash, and issued early on the third day, and hired two men, who
        put it right." "It isn't strange then!" exclaimed Pao-yü, "When the lotus
        blossomed last month in the pond of our garden of Broad Vista, I plucked ten
        of them and bade T'sai Ming go out of town and lay them as my offering on
        his grave. On his return, I also inquired of him: whether it had been
        damaged by the water or not; and he explained that not only had it not
        sustained any harm, but that it looked better than when last he'd seen it.}
}
{
}

\Verse{54}
{
    \zA{}
}
{
    \eA{Several of his friends, I argued, must have had it put in proper repair; and
        I felt it irksome that I should, day after day, be so caged at home as to be
        unable to be my own master in the least thing, and that if even I move, and
        any one comes to know of it, this one is sure to exhort me, if that one does
        not restrain me. I can thus afford to brag, but can't manage to act! And
        though I've got plenty of money, I'm not at liberty to spend any of it!"
        "There's no use your worrying in a matter like this!" Liu Hsiang-lien said.
        "I am outside, so all you need do is to inwardly foster the wish; that's
        all. But as the first of the tenth moon will shortly be upon us, I've
        already prepared the money necessary for going to the graves. You know well
        enough that I'm as poor as a rat;}
}
{
}

\Verse{55}
{
    \zA{}
}
{
    \eA{I've no hoardings at home; and when a few cash find their way into my
        pocket, I soon remain again quite empty-handed. But I'd better make the best
        of this opportunity, and keep the amount I have, in order that, when the
        time comes, I mayn't find myself without a cash." "It's exactly about this
        that I meant to send Pei Ming to see you," Pao-yü added. "But it isn't often
        that one can manage to find you at home. I'm well aware how uncertain your
        movements are; one day you are here, and another there; you've got no fixed
        resort." "There's no need sending any one to hunt me up!" Liu Hsiang-lien
        replied. "All that each of us need do in this matter is to acquit ourselves
        of what's right. But in a little while, I again purpose going away on a tour
        abroad, to return in three to five years' time."}
}
{
}

\Verse{56}
{
    \zA{}
}
{
    \eA{When Pao-yü heard his intention, "Why is this?" he at once inquired. Liu
        Hsiang-lien gave a sardonic smile. "When my wish is on a fair way to be
        accomplished," he said, "you'll certainly hear everything. I must now leave
        you." "After all the difficulty we've had in meeting," Pao-yü remarked,
        "wouldn't it be better were you and I to go away together in the evening?"
        "That worthy cousin of yours," Hsiang-lien rejoined, "is as bad as ever, and
        were I to stay any longer, trouble would inevitably arise. So it's as well
        that I should clear out of his way." Pao-yü communed with himself for a
        time. "In that case," he then observed, "it's only right, that you should
        retire. But if you really be bent upon going on a distant tour, you must
        absolutely tell me something beforehand.}
}
{
}

\Verse{57}
{
    \zA{}
}
{
    \eA{Don't, on any account, sneak away quietly!". As he spoke, the tears trickled
        down his cheeks. "I shall, of course, say good-bye to you," Liu Hsiang-lien
        rejoined. "But you must not let any one know anything about it!" While
        uttering these words, he stood up to get away. "Go in at once," he urged,
        "there's no need to see me off!" Saying this, he quitted the study. As soon
        as he reached the main entrance, he came across Hsüeh P'an, bawling out
        boisterously, "Who let young Liu-erh go?" The moment these shouts fell on
        Liu Hsiang-lien's ear, his anger flared up as if it had been sparks spurting
        wildly about, and he only wished he could strike him dead with one blow.}
}
{
}

\Verse{58}
{
    \zA{}
}
{
    \eA{But on second consideration, he pondered that a fight after the present
        festive occasion would be an insult to Lai Shang-jung, and he perforce felt
        bound to stifle his indignation. When Hsüeh P'an suddenly espied him walking
        out, he looked as delighted as if he had come in for some precious gem. With
        staggering step he drew near him. Clutching him with one grip, "My dear
        brother," he smirked. "where are you off to?" "I'm going somewhere, but will
        be back soon," Hsiang-lien said by way of response. "As soon as you left,"
        Hsüeh P'an smiled, "all the fun went. But pray sit a while! If you do so, it
        will be a proof of your regard for me! Don't flurry yourself. With such a
        senior brother as myself to stand by you, it will be as easy a job for you
        to become an official as to reap a fortune."}
}
{
}

\Verse{59}
{
    \zA{}
}
{
    \eA{The sight of his repulsive manner filled the heart of Hsiang-lien with
        disgust and shame. But speedily devising a plan, he drew him to a secluded
        spot. "Is your friendship real," he smiled, "or is it only a sham?" This
        question sent Hsüeh P'an into such raptures that he found it difficult to
        check himself from gratifying his longings. But glancing at him with the
        corner of his eye, "My dear brother," he smiled, "what makes you ask me such
        a thing? If my friendship for you is a sham, may I die this moment, before
        your very eyes." "Well, if that be so," Hsiang-lien proceeded, "it isn't
        convenient in here, so sit down and wait a bit.}
}
{
}

\Verse{60}
{
    \zA{}
}
{
    \eA{I'll go ahead, but come out of this yourself by and bye, and follow me to my
        place, where we can drink the whole night long. I've also got there two
        first-rate young fellows who never go out of doors. But don't bring so much
        as a single follower with you, as you'll find, when you get there, plenty of
        people ready at hand to wait on you." So high did this assignation raise
        Hsüeh P'an's spirits that he recovered, to a certain extent, from the
        effects of wine. "Is it really so?" he asked. "How is it," Hsiang-lien
        laughed, "that when people treat you with a sincere heart, you don't, after
        all, believe them?" "I'm no fool," eagerly exclaimed Hsüeh P'an, "and how
        could I not believe you?}
}
{
}

\Verse{61}
{
    \zA{}
}
{
    \eA{But since this be the case, how am I, who don't even know the way, to find
        your whereabouts if you are to go ahead of me?" "My place is outside the
        northern gate." Hsiang-lien explained. "But can you tear yourself away from
        your home to spend the night outside the city walls?" "As long as you're
        there," Hsüeh P'an said, "what will I want my home for?" "If that be so,"
        Hsiang-lien resumed, "I'll wait for you on the bridge outside the northern
        gate. But let us meanwhile rejoin the banquet and have some wine. Come
        along, after you've seen me go; they won't notice us then." "Yes!" shouted
        Hsüeh P'an with alacrity as he acquiesced to the proposal. The two young
        fellows thereupon returned to the feast, and drank for a time.}
}
{
}

\Verse{62}
{
    \zA{}
}
{
    \eA{Hsüeh Pan, however, could with difficulty endure the suspense. He kept his
        gaze intent upon Hsiang-lien; and the more he pondered within himself upon
        what was coming, the more exuberance swelled in his heart. Now he emptied
        one wine-kettle; now another; and, without waiting for any one to press him,
        he, of his own accord, gulped down one drink after another, with the result
        that he unconsciously made himself nearly quite tipsy. Hsiang-lien then got
        up and quitted the room, and perceiving every one off his guard, he egressed
        out of the main entrance. "Go home ahead," he directed his page Hsing Nu.
        "I'm going out of town, but I'll be back at once."}
}
{
}

\Verse{63}
{
    \zA{}
}
{
    \eA{By the time he had finished giving him these directions, he had already
        mounted his horse, and straightway he proceeded to the bridge beyond the
        northern gate, and waited for Hsüeh P'an. A long while elapsed, however,
        before he espied Hsüeh P'an in the distance, hurrying along astride of a
        high steed, with gaping mouth, staring eyes, and his head, banging from side
        to side like a pedlar's drum. Without intermission, he glanced confusedly
        about, sometimes to the left, and sometimes to the right; but, as soon as he
        got where he had to pass in front of Hsiang-lien's horse, he kept his gaze
        fixed far away, and never troubled his mind with the immediate vicinity.}
}
{
}

\Verse{64}
{
    \zA{}
}
{
    \eA{Hsiang-lien felt amused and angry with him, but forthwith giving his horse
        also the rein, he followed in his track, while Hsüeh P'an continued to stare
        ahead. Little by little the habitations got scantier and scantier, so
        pulling his horse round, (Hsüeh P'an) retraced his steps. The moment he
        turned back, he unawares caught sight of Hsiang-lien, and his spirits rose
        within him, as if he had got hold of some precious thing of an extraordinary
        value. "I knew well enough," he eagerly smiled, "that you weren't one to
        break faith." "Quick, let's go ahead!" Hsiang-lien smilingly urged. "Mind
        people might notice us and follow us. It won't then be nice!"}
}
{
}

\Verse{65}
{
    \zA{}
}
{
    \eA{While instigating him, he took the lead, and letting his horse have the
        rein, he wended his way onwards, followed closely by Hsüeh P'an. But when
        Hsiang-lien perceived that the country ahead of them was already thinly
        settled and saw besides a stretch of water covered with a growth of weeds,
        he speedily dismounted, and tied his horse to a tree. Turning then round;
        "Get down!" he said, laughingly, to Hsüeh P'an. "You must first take an
        oath, so that in the event of your changing your mind in the future, and
        telling anything to anyone, the oath might be accomplished." "You're quite
        right!" Hsüeh P'an smiled; and jumping down with all despatch, he too made
        his horse fast to a tree, and then crouched on his knees.}
}
{
}

\Verse{66}
{
    \zA{}
}
{
    \eA{"If I ever in days to come," he exclaimed, "know any change in my feelings
        and breathe a word to any living soul, may heaven blast me and earth
        annihilate me!" Scarcely had he ended this oath, when a crash fell on his
        ear, and lo, he felt as if an iron hammer had been brought down to bear upon
        him from behind. A black mist shrouded his eyes, golden stars flew wildly
        about before his gaze; and losing all control over himself, he sprawled on
        the ground. Hsiang-lien approached and had a look at him; and, knowing how
        little he was accustomed to thrashings, he only exerted but little of his
        strength, and struck him a few blows on the face.}
}
{
}

\Verse{67}
{
    \zA{}
}
{
    \eA{But about this time a fruit shop happened to open, and Hsüeh P'an strained
        at first every nerve to rise to his feet, when another slight kick from
        Hsiang-lien tumbled him over again. "Both parties should really be
        agreeable," he shouted. "But if you were not disposed to accept my advances,
        you should have simply told me in a proper way. And why did you beguile me
        here to give me a beating?" So speaking, he went on boisterously to heap
        invective upon his head. "I'll take you, you blind fellow, and show you who
        Mr. Liu is," Hsiang-lien cried. "You don't appeal to me with solicitous
        entreaties, but go on abusing me! To kill you would be of no use, so I'll
        merely give you a good lesson!" With these words, he fetched his whip, and
        administered him, thirty or forty blows from his back down to his shins.}
}
{
}

\Verse{68}
{
    \zA{}
}
{
    \eA{Hsüeh P'an had sobered down considerably from the effects of wine, and found
        the stings of pain so intolerable, that little able to restrain himself, he
        gave way to groans. "Do you go on in this way?" Hsiang-lien said, with an
        ironical smile. "Why, I thought you were not afraid of beatings." While
        uttering this taunt, he seized Hsüeh P'an by the left leg, and dragging him
        several steps into a miry spot among the reeds, he rolled him about till he
        was covered with one mass of mud. "Do you now know what stuff I'm made of?"
        he proceeded to ask. Hsüeh P'an made no reply. But simply lay prostrate, and
        moaned. Then throwing away his whip Hsiang-lien gave him with his fist
        several thumps all over the body.}
}
{
}

\Verse{69}
{
    \zA{}
}
{
    \eA{Hsüeh P'an began to wriggle violently and vociferate wildly. "Oh, my ribs
        are broken!" he shouted. "I know you're a proper sort of person! It's all
        because I made the mistake of listening to other people's gossip!" "There's
        no need for you to drag in other people!" Hsiang-lien went on. "Just confine
        yourself to those present!" "There's nothing up at present!" Hsüeh P'an
        cried. "From what you say, you're a person full of propriety. So it's I who
        am at fault." "You'll have to speak a little milder," Hsiang-lien added,
        "before I let you off." "My dear younger brother," Hsüeh P'an pleaded, with
        a groan. Hsiang-lien at this struck him another blow with his fist. "Ai!"
        ejaculated Hsüeh P'an. "My dear senior brother!" he exclaimed. Hsiang-lien
        then gave him two more whacks, one after the other.}
}
{
}

\Verse{70}
{
    \zA{}
}
{
    \eA{"Ai Yo!" Hsüeh P'an precipitately screamed. "My dear Sir, do spare me, an
        eyeless beggar; and henceforth I'll look up to you with veneration; I'll
        fear you!" "Drink two mouthfuls of that water!" shouted Hsiang-lien. "That
        water is really too foul," Hsüeh P'an argued, in reply to this suggestion,
        wrinkling his eyebrows the while; "and how could I put any of it in my
        mouth?" Hsiang-lien raised his fist and struck him. "I'll drink it, I'll
        drink it!" quickly bawled Hsüeh P'an. So saying, he felt obliged to lower
        his head to the very roots of the reeds and drink a mouthful. Before he had
        had time to swallow it, a sound of 'ai' became audible, and up came all the
        stuff he had put into his mouth only a few seconds back. "You filthy thing!"
        exclaimed Hsiang-lien. "Be quick and finish drinking; and I'll let you off."}
}
{
}

\Verse{71}
{
    \zA{}
}
{
    \eA{Upon hearing this, Hsüeh P'an bumped his head repeatedly on the ground. "Do
        please," he cried, "lay up a store of meritorious acts for yourself and let
        me off! I couldn't take that were I even on the verge of death!" "This kind
        of stench will suffocate me!" Hsiang-lien observed, and, with this remark,
        he abandoned Hsüeh Pan to his own devices; and, pulling his horse, he put
        his foot to the stirrup, and rode away. Hsüeh Pan, meanwhile, became aware
        of his departure, and felt at last relieved in his mind. Yet his conscience
        pricked him for he saw that he should not misjudge people. He then made an
        effort to raise himself, but the racking torture he experienced all over his
        limbs was so sharp that he could with difficulty bear it.}
}
{
}

\Verse{72}
{
    \zA{}
}
{
    \eA{Chia Chen and the other guests present at the banquet became, as it
        happened, suddenly alive to the fact that the two young fellows had
        disappeared; but though they extended their search everywhere, they saw
        nothing of them. Some one insinuated, in an uncertain way, that they had
        gone outside the northern gate; but as Hsüeh P'an's pages had ever lived in
        dread of him, who of them had the audacity to go and hunt him up after the
        injunctions, he had given them, that they were not to follow him? But waxing
        solicitous on his account, Chia Chen subsequently bade Chia Jung take a few
        servant-boys and go and discover some clue of him, or institute inquiries as
        to his whereabouts.}
}
{
}

\Verse{73}
{
    \zA{}
}
{
    \eA{Straightway therefore they prosecuted their search beyond the northern gate,
        to a distance of two li below the bridge, and it was quite by accident that
        they discerned Hsüeh P'an's horse made fast by the side of a pit full of
        reeds. "That's a good sign!" they with one voice exclaimed; "for if the
        horse is there, the master must be there too!" In a body, they thronged
        round the horse, when, from among the reeds, they caught the sound of human
        groans, so hurriedly rushing forward to ascertain for themselves, they, at a
        glance, perceived Hsüeh P'an, his costume all in tatters, his countenance
        and eyes so swollen and bruised that it was hard to make out the head and
        face, and his whole person, inside as well as outside his clothes, rolled
        like a sow in a heap of mud.}
}
{
}

\Verse{74}
{
    \zA{}
}
{
    \eA{Chia Jung surmised pretty nearly the truth. Speedily dismounting, he told
        the servants to prop him up. "Uncle Hsüeh," he laughed, "you daily go in for
        lewd dalliance; but have you to-day come to dissipate in a reed-covered pit?
        The King of the dragons in this pit must have also fallen in love with your
        charms, and enticed you to become his son-in-law that you've come and gored
        yourself on his horns like this!" Hsüeh P'an was such a prey to intense
        shame that he would fain have grovelled into some fissure in the earth had
        he been able to detect any. But so little able was he to get on his horse
        that Chia Jung directed a servant to run to the suburbs and fetch a chair.
        Ensconced in this, Hsüeh P'an entered town along with the search party.}
}
{
}

\Verse{75}
{
    \zA{}
}
{
    \eA{Chia Jung still insisted upon carrying him to Lai Ta's house to join the
        feast, so Hsüeh P'an had to make a hundred and one urgent appeals to him to
        tell no one, before Chia Jung eventually yielded to his solicitations and
        allowed him to have his own way and return home. Chia Jung betook himself
        again to Lai Ta's house, and narrated to Chia Chen their recent experiences.
        When Chia Chen also learnt of the flogging (Hsüeh P'an) had received from
        Hsiang-lien, he laughed. "It's only through scrapes," he cried, "that he'll
        get all right!" In the evening, after the party broke up, he came to inquire
        after him. But Hsüeh P'an, who was lying all alone in his bedroom, nursing
        himself, refused to see him, on the plea of indisposition.}
}
{
}

\Verse{76}
{
    \zA{}
}
{
    \eA{When dowager lady Chia and the other inmates had returned home, and every
        one had retired into their respective apartments, Mrs. Hsüeh and Pao-ch'ai
        observed that Hsiang Ling's eyes were quite swollen from crying, and they
        questioned her as to the reason of her distress. (On being told), they
        hastily rushed to look up Hsüeh P'an; but, though they saw his body covered
        with scars, they could discover no ribs broken, or bones dislocated. Mrs.
        Hsüeh fell a prey to anguish and displeasure. At one time, she scolded Hsüeh
        P'an; at another, she abused Liu Hsiang-lien. Her wish was to lay the matter
        before Madame Wang in order that some one should be despatched to trace Liu
        Hsiang-lien and bring him back, but Pao-ch'ai speedily dissuaded her.}
}
{
}

\Verse{77}
{
    \zA{}
}
{
    \eA{"It's nothing to make a fuss about," she represented. "They were simply
        drinking together; and quarrels after a wine bout are ordinary things. And
        for one who's drunk to get a few whacks more or less is nothing uncommon!
        Besides, there's in our home neither regard for God nor discipline. Every
        one knows it. If it's purely out of love, mother, that you desire to give
        vent to your spite, it's an easy matter enough. Have a little patience for
        three or five days, until brother is all right and can go out. Mr. Chia Chen
        and Mr. Chia Lien over there are not people likely to let the affair drop
        without doing anything! They'll, for a certainty, stand a treat, and ask
        that fellow, and make him apologise and admit his wrong in the presence of
        the whole company, so that everything will be properly settled.}
}
{
}

\Verse{78}
{
    \zA{}
}
{
    \eA{But were you now, ma, to begin making much of this occurrence, and telling
        every one, it would, on the contrary, look as if you had, in your motherly
        partiality and fond love for him, indulged him to stir up a row and provoke
        people! He has, on this occasion, had unawares to eat humble pie, but will
        you, ma, put people to all this trouble and inconvenience and make use of
        the prestige enjoyed by your relatives to oppress an ordinary person?" "My
        dear child," Mrs. Hsüeh rejoined, "after listening to the advice proffered
        by her, you've, after all, been able to foresee all these things! As for me,
        that sudden fit of anger quite dazed me!"}
}
{
}

\Verse{79}
{
    \zA{}
}
{
    \eA{"All will thus be square," Pao-ch'ai smiled, "for, as he's neither afraid of
        you, mother, nor gives an ear to people's exhortations, but gets wilder and
        wilder every day that goes by, he may, if he gets two or three lessons, turn
        over a new leaf." While Hsüeh P'an lay on the stovecouch, he reviled
        Hsiang-lien with all his might. Next, he instigated the servant-boys to go
        and demolish his house, kill him and bring a charge against him.}
}
{
}

\Verse{80}
{
    \zA{}
}
{
    \eA{But Mrs. Hsüeh hindered the lads from carrying out his purpose, and
        explained to her son: "that Liu Hsiang-lien had casually, after drinking,
        behaved in a disorderly way, that now that he was over the effects of wine,
        he was exceedingly filled with remorse, and that, prompted by the fear of
        punishment, he had effected his escape." But, reader, if you feel any
        interest to know what happened when Hsüeh P'an heard the version his mother
        gave him, listen to what you will find in the next chapter.}
}
{
}
